\section{Журнал выполнения}
\begin{enumerate}
    \item На машине нет \texttt{matplotlib} и \texttt{gnuplot}. Чтобы не добавлять зависимостей, график строится пакетом \texttt{pgfplots} прямо в отчёте, а скрипт только готовит файл точек.
    \item Строить все 3,97 млн точек не нужно: \texttt{pgfplots} не хватает памяти \TeX, а в логарифмическом масштабе точки хвоста сливаются. Поэтому ранги прорежены с постоянным шагом в логарифмической шкале (через 3\,\%), форма кривой, в т.\,ч. ступеньки хвоста, при этом сохраняется.
    \item Все файлы корпуса подаются одному процессу токенизатора: запуск процесса на каждый из 1\,000\,000 файлов занял бы больше времени, чем сам подсчёт. Чтобы последний токен файла не склеился с первым токеном следующего, между файлами вставляется перевод строки (тест 3). Подсчёт частот по всему корпусу (995 млн токенов) занимает около 5~мин, из них больше половины --- чтение миллиона файлов.
    \item У терминов с одинаковой частотой порядок (и значит ранг) условный --- по алфавиту. Например, термин на ранге 3\,974\,421 --- просто последний по алфавиту среди 2\,042\,771 термина с частотой 1. Символы, которых нет в шрифте отчёта, в таблице записаны как \texttt{[U+XXXX]}.
\end{enumerate}

\section{План тестирования}
Тесты (\texttt{tests/run\_tests.py}) запускают \texttt{scripts/zipf.py} как утилиту командной строки на маленьких корпусах, созданных во временном каталоге, и проверяют выходные файлы.

\begin{center}
\small
\begin{tabular}{|c|p{0.3\textwidth}|p{0.58\textwidth}|}
\hline
\textnumero & Тест & Что проверяется \\
\hline
1 & Частоты и ранги & Два файла во вложенных каталогах (и файл \texttt{.md}, который не должен читаться): \texttt{freq.tsv} совпадает с посчитанным вручную --- порядок по убыванию частоты, равные частоты по алфавиту, регистр понижен; число токенов, терминов, терминов с частотой 1 и $C = f(1)$. \\
\hline
2 & Точки графика & Корпус из 200 слов с частотами $\lfloor 200/r \rfloor$: в \texttt{zipf.dat} первая точка --- ранг 1, последняя --- ранг 200 (число терминов), ранги строго возрастают, частота в каждой точке совпадает с заданной, столбец закона равен $C/r$. \\
\hline
3 & Граница файлов & Файлы \texttt{foo} и \texttt{bar} без перевода строки в конце: термины \texttt{foo} и \texttt{bar}, а не \texttt{foobar}. \\
\hline
4 & Ошибки & Пустой каталог и корпус без токенов (только знаки препинания): ненулевой код возврата и сообщение в \texttt{stderr}. \\
\hline
\end{tabular}
\end{center}

\section{Результаты}
\subsection{Автотесты}
Пройдены все 4 теста. Журнал прогона (\texttt{results/tests.log}):

{\scriptsize
\VerbatimInput[breaklines, breakanywhere]{../results/tests.log}
}

\subsection{Распределение терминов по частотам}
Корпус: 1\,000\,000 документов, 995\,461\,170 токенов, 3\,974\,421 разный термин, из них 2\,042\,771 (51,4\,\%) встречаются ровно один раз. Самый частый термин --- \texttt{the}, $C = f(1) = 70\,066\,836$.

\begin{center}
\begin{tikzpicture}
\begin{loglogaxis}[
    width=0.95\textwidth, height=0.72\textwidth,
    xlabel={Ранг термина $r$}, ylabel={Частота $f$},
    xmin=1, xmax=1e7, ymin=0.5, ymax=2e8,
    grid=major, legend pos=north east, legend cell align=left]
\addplot[blue, thick, mark=*, mark size=0.8pt] table[x=rank, y=freq] {../results/zipf.dat};
\addplot[red, thick, dashed] table[x=rank, y=zipf] {../results/zipf.dat};
\legend{Корпус, {Закон Ципфа $f = C/r$, $C = 70\,066\,836$}}
\end{loglogaxis}
\end{tikzpicture}
\end{center}

Частоты на отдельных рангах (\texttt{results/freq.tsv}); последний столбец --- во сколько раз реальная частота больше предсказанной законом:
\begin{center}
\footnotesize
\begin{tabular}{|r|l|r|r|r|r|}
\hline
$r$ & Термин & $f$ & $f \cdot r$ & $C/r$ & $f / (C/r)$ \\
\hline
1 & \texttt{the} & 70\,066\,836 & 70\,066\,836 & 70\,066\,836,0 & 1,00 \\
\hline
2 & \texttt{of} & 32\,227\,461 & 64\,454\,922 & 35\,033\,418,0 & 0,92 \\
\hline
3 & \texttt{and} & 29\,184\,703 & 87\,554\,109 & 23\,355\,612,0 & 1,25 \\
\hline
10 & \texttt{on} & 7\,945\,476 & 79\,454\,760 & 7\,006\,683,6 & 1,13 \\
\hline
100 & \texttt{station} & 792\,854 & 79\,285\,400 & 700\,668,4 & 1,13 \\
\hline
500 & \texttt{private} & 216\,889 & 108\,444\,500 & 140\,133,7 & 1,55 \\
\hline
1\,000 & \texttt{match} & 117\,605 & 117\,605\,000 & 70\,066,8 & 1,68 \\
\hline
10\,000 & \texttt{aforementioned} & 7\,188 & 71\,880\,000 & 7\,006,7 & 1,03 \\
\hline
100\,000 & \texttt{homestyle} & 168 & 16\,800\,000 & 700,7 & 0,24 \\
\hline
1\,000\,000 & \texttt{edwill} & 4 & 4\,000\,000 & 70,1 & 0,06 \\
\hline
1\,931\,651 & \texttt{0000000000} & 1 & 1\,931\,651 & 36,3 & 0,03 \\
\hline
3\,974\,421 & \texttt{[U+313BC][U+5BB6][U+4EBA]} & 1 & 3\,974\,421 & 17,6 & 0,06 \\
\hline
\end{tabular}

\end{center}

\subsection{Причины расхождения}
Кривая корпуса похожа на прямую закона только в целом: частота падает примерно степенным образом на семь порядков рангов. Но наклон у неё другой и меняется по длине, и можно выделить три участка.

\textbf{Начало ($r \le 70$).} Здесь закон выполняется лучше всего: реальная частота отличается от предсказанной не больше чем в 1,25 раза. Но частоты самых частых слов убывают не строго как $1/r$: \texttt{of} и \texttt{and} (ранги 2 и 3) почти равны по частоте (32,2 и 29,2 млн), хотя по закону должны отличаться в 1,5 раза. Частоты служебных слов (\texttt{the}, \texttt{of}, \texttt{and}, \texttt{in}, \texttt{to}, \texttt{a}) определяет грамматика английского языка, а не единая закономерность для всех слов. Именно отклонения начала кривой учитывает параметр $\rho$ в формуле Мандельброта $f = P / (r + \rho)^B$.

\textbf{Середина ($70 < r \le 10\,417$).} Кривая идёт выше прямой, до 1,73 раза около ранга 1\,769. От ранга 10 до 1000 частота падает в 68 раз (наклон в логарифмических осях $-0{,}92$), а по закону должна в 100 раз. Средние ранги \enquote{переполнены} терминами из-за особенностей корпуса:
\begin{itemize}
    \item корпус тематический: слова туризма встречаются гораздо чаще, чем в общем английском тексте (\texttt{city} --- ранг 45, \texttt{museum} --- 177, \texttt{travel} --- 210, \texttt{hotel} --- 245, \texttt{trip} --- 262, \texttt{beach} --- 423);
    \item корпус смешан из источников с разной лексикой: энциклопедия (Wikipedia, 62\,\% документов), путеводитель (Wikivoyage) и разговорные посты и ответы (Reddit, Travel SE). Из постов --- местоимения первого и второго лица: \texttt{i} --- ранг 20, \texttt{you} --- 58;
    \item числа: 73 из первой тысячи терминов --- числа (годы, \texttt{00} из времени \enquote{10:00} --- ранг 415, цены), это не слова языка;
    \item осколки слов, которые даёт токенизатор ЛР3: \texttt{s} из \texttt{'s} (ранг 13), \texttt{t} из \texttt{don't} (119), буквы сторон света \texttt{e}, \texttt{w}, \texttt{n} из координат (ранги 389, 932, 1\,469); части адресов сайтов \texttt{com}, \texttt{https}, \texttt{www}.
\end{itemize}

\textbf{Хвост ($r > 10\,417$).} Кривая уходит под прямую, ниже всего --- в 36 раз (0,028) в начале \enquote{полки} частоты 1 (ранг 1\,931\,651). От ранга $10^4$ до $10^5$ частота падает в 43 раза (наклон $-1{,}63$), от $10^5$ до $10^6$ --- в 42 раза, а по закону должна в 10 раз. Причины:
\begin{itemize}
    \item конечный размер корпуса. Закон с $C = 70\,066\,836$ предсказывает около 70 млн разных терминов (частота 1 на ранге $r = C$), а их 3,97 млн: редкие слова, которые появились бы в корпусе большего объёма, в этом просто не встретились;
    \item частота --- целое число, меньше 1 не бывает. Поэтому хвост --- лестница горизонтальных ступенек ($f = 3, 2, 1$), а последняя ступенька --- 2\,042\,771 термин с частотой 1 на рангах от 1\,931\,651 до 3\,974\,421 --- горизонтальная полка. На ней кривая снова приближается к прямой (0,06 на последнем ранге), но только потому, что прямая продолжает опускаться;
    \item хвост состоит в основном не из слов языка: среди терминов с частотой 1 у 12,5\,\% есть цифры (ISBN, телефоны, цены), у 18,4\,\% --- не-ASCII символы (иностранные слова, осколки тайских слов, фразы на языках без пробелов до 75 символов), остальное --- редкие имена собственные, опечатки, склеенные слова и куски ссылок из постов (самый длинный токен --- 754 символа base64 из адреса ссылки).
\end{itemize}

\textbf{Сравнение с корпусом в 30\,000 документов.} На первом корпусе ЛР1 (36,7 млн токенов) хвост был ниже прямой не больше чем в 6 раз, а на миллионе --- в 36 раз. Константа $C = f(1)$ растёт вместе с числом токенов (токенов стало больше в 27 раз, $C$ --- в 33 раза: 2,1 $\to$ 70 млн), а словарь --- намного медленнее, по закону Хипса (в 5 раз: 775 тыс. $\to$ 3,97 млн). Поэтому чем больше корпус, тем дальше от прямой с $C = f(1)$ уходит хвост; доля терминов с частотой 1 при этом снизилась с 57 до 51\,\%.

В целом закон Ципфа --- эмпирическое приближение, которое, как и сказано в лекции, хорошо описывает только средние по частоте термины и ошибается в начале и в хвосте словаря. Кроме того, у нашего корпуса нет нормализации словоформ (\texttt{museum}/\texttt{museums} --- разные термины) и много нетекстовых токенов, что сдвигает распределение относительно \enquote{чистого} языка.

\section{Выводы}
Построен частотный список терминов корпуса из 1\,000\,000 документов (995 млн токенов, 3\,974\,421 термин) и график распределения в логарифмической шкале с наложенным законом Ципфа $f = C/r$, $C = f(1)$. Кривая корпуса в целом следует степенному закону, но отклоняется от него: в середине (ранги 70--10\,417) частоты выше предсказанных до 1,7 раза из-за тематической лексики, смеси источников, чисел и осколков слов, а в хвосте --- до 36 раз ниже из-за конечного размера корпуса и целочисленных частот (51\,\% терминов встречаются один раз). С ростом корпуса хвост отклоняется сильнее: словарь растёт медленнее числа токенов. Точнее распределение описала бы формула Мандельброта с подобранными константами (дополнительное задание, не выполнялось). Ближе к закону распределение стало бы после очистки корпуса от чисел и мусорных токенов и нормализации словоформ (ЛР4).
